\documentclass[11pt, a4paper]{article}
\usepackage[a4paper, margin=2.5cm]{geometry}
\usepackage{fontspec}
\usepackage{enumitem}
\usepackage{xcolor}
\usepackage{listings}
\usepackage{tcolorbox}
\usepackage[hidelinks]{hyperref}

% Definicao de cores do documento
\definecolor{codebg}{RGB}{245, 247, 250}
\definecolor{codeframe}{RGB}{210, 215, 225}
\definecolor{primary}{RGB}{15, 76, 129}
\definecolor{accent}{RGB}{40, 116, 166}
\definecolor{warnbg}{RGB}{253, 237, 237}
\definecolor{warnframe}{RGB}{220, 53, 69}
\definecolor{infobg}{RGB}{240, 248, 255}
\definecolor{infoframe}{RGB}{40, 116, 166}

% Linguagem YAML personalizada para listings
\lstdefinelanguage{yaml}{
  keywords={true, false, null, y, n},
  keywordstyle=\color{primary}\bfseries,
  basicstyle=\ttfamily\small,
  sensitive=false,
  comment=[l]{\#},
  commentstyle=\color{gray}\itshape,
  stringstyle=\color{accent},
  morestring=[b]',
  morestring=[b]"
}

% Configuracao basica do listings
\lstset{
  backgroundcolor=\color{codebg},
  frame=single,
  rulecolor=\color{codeframe},
  basicstyle=\ttfamily\small,
  breaklines=true,
  columns=fullflexible,
  keepspaces=true,
  showstringspaces=false,
  commentstyle=\color{gray}
}

\setlength{\parindent}{0pt}
\setlength{\parskip}{7pt}

\begin{document}

\begin{center}
    {\LARGE \textbf{\color{primary} Laboratório de Kubernetes}} \\[4pt]
    {\large \textbf{\color{accent} Da Máquina Local à Nuvem Híbrida (Windows + kind)}} \\[8pt]
    \rule{\textwidth}{1pt}
\end{center}

\textbf{Objetivo:} Sair do zero, criar um cluster Kubernetes local multi-nó, testar a auto-recuperação e, por fim, simular uma arquitetura de nuvem distribuída.

\begin{tcolorbox}[colback=infobg, colframe=infoframe, arc=3pt, boxrule=0.8pt]
\textbf{INSTRUÇÕES DE ENTREGA DA ATIVIDADE:}
\begin{enumerate}[leftmargin=*, topsep=3pt, itemsep=2pt]
    \item Execute todos os comandos detalhados neste roteiro prático.
    \item Capture telas (\textit{prints}) que comprovem a execução e os resultados dos comandos no seu terminal.
    \item Elabore e entregue um \textbf{relatório individual} explicando o que foi feito em cada etapa, ilustrado com os seus respectivos \textit{prints}.
\end{enumerate}
\end{tcolorbox}

\begin{tcolorbox}[colback=warnbg, colframe=warnframe, arc=3pt, boxrule=0.8pt]
\textbf{PRÉ-REQUISITOS (WINDOWS):}
\begin{itemize}[leftmargin=*, topsep=3pt, itemsep=2pt]
    \item Docker Desktop instalado e rodando (com WSL2 ativado).
    \item Terminal PowerShell aberto como Administrador.
    \item Instale o \texttt{kubectl} e o \texttt{kind} via Chocolatey:
    \begin{lstlisting}[language=bash]
choco install kubernetes-cli kind -y
    \end{lstlisting}
\end{itemize}
\end{tcolorbox}

\section*{Fase 1: O Básico (Comandos e Deploy)}

Vamos criar o nosso "datacenter" local com 1 nó de controle e 2 servidores de trabalho usando o \texttt{kind}.

\subsection*{1. Criar o cluster multi-nó}
Crie um arquivo chamado \texttt{kind-config.yaml} na sua pasta de trabalho com o seguinte conteúdo (ele mapeia a porta 30080 para que você consiga acessar no navegador):

\begin{lstlisting}[language=yaml]
kind: Cluster
apiVersion: kind.x-k8s.io/v1alpha4
nodes:
- role: control-plane
  extraPortMappings:
  - containerPort: 30080
    hostPort: 30080
    protocol: TCP
- role: worker
- role: worker
\end{lstlisting}

No PowerShell, crie o cluster:
\begin{lstlisting}[language=bash]
kind create cluster --name lab-k8s --config kind-config.yaml
\end{lstlisting}

Verifique se os 3 nós estão prontos:
\begin{lstlisting}[language=bash]
kubectl get nodes
\end{lstlisting}

\subsection*{2. Criar um Deployment (Modo Declarativo)}
Vamos criar um gerenciador (\texttt{Deployment}) que garante 3 réplicas do Nginx:

\begin{lstlisting}[language=bash]
kubectl create deployment meu-app --image=nginx:latest --replicas=3
\end{lstlisting}

\subsection*{3. Expor e Testar}
Exponha a aplicação na porta mapeada:

\begin{lstlisting}[language=bash]
kubectl expose deployment meu-app --type=NodePort --port=80 --target-port=80 --node-port=30080
\end{lstlisting}

\textit{Acesse no navegador: \texttt{http://localhost:30080}. Você verá a página de boas-vindas do Nginx.}

\section*{Fase 2: Testando Conceito Local (Auto-cura)}

Vamos provar que o Deployment mantém o estado desejado.

\textbf{1. Simule uma falha:} \\
Pegue o nome de um dos pods listados abaixo (será algo como \texttt{meu-app-xxxxxxxx-xxxxx}):

\begin{lstlisting}[language=bash]
kubectl get pods
\end{lstlisting}

Delete o pod (substitua pelo nome que você copiou):

\begin{lstlisting}[language=bash]
kubectl delete pod meu-app-<hash>
\end{lstlisting}

\textbf{2. Observando a mágica:} \\
Rode o comando com \texttt{-w} (\textit{watch}) para ver em tempo real:

\begin{lstlisting}[language=bash]
kubectl get pods -w
\end{lstlisting}

\begin{tcolorbox}[colback=infobg, colframe=infoframe, arc=3pt, boxrule=0.8pt]
\textbf{OBSERVE:} O pod que você deletou entrou em \texttt{Terminating} e um \textbf{novo pod foi criado automaticamente} em segundos. O Kubernetes percebeu que você queria 3 e recriou o que faltava. \textit{(Aperte \texttt{Ctrl+C} para sair do watch).}
\end{tcolorbox}

\subsection*{3. LIMPEZA (A Transição)}
Agora que já entendemos o básico, vamos apagar o \texttt{meu-app} para que a nossa próxima aplicação seja a única a rodar no cluster.

\begin{lstlisting}[language=bash]
kubectl delete deployment meu-app
kubectl delete service meu-app
\end{lstlisting}

\section*{Fase 3: A Nuvem Híbrida (Distribuição entre Servidores)}

Agora, vamos simular um ambiente de nuvem distribuída. O seu cluster possui 3 máquinas virtuais (nós). Vamos forçar o Kubernetes a colocar uma réplica em cada nó diferente para garantir resiliência.

\subsection*{1. Deploy Distribuído (O YAML Protegido)}
Crie um arquivo chamado \texttt{app-nuvem.yaml} com o código abaixo. Ele já está 100\% configurado.

\begin{lstlisting}[language=yaml]
apiVersion: apps/v1
kind: Deployment
metadata:
  name: app-distribuida
spec:
  replicas: 2
  selector:
    matchLabels:
      app: distribuida
  template:
    metadata:
      labels:
        app: distribuida
    spec:
      affinity:
        podAntiAffinity:
          requiredDuringSchedulingIgnoredDuringExecution:
            - labelSelector:
                matchExpressions:
                  - key: app
                    operator: In
                    values:
                      - distribuida
              topologyKey: kubernetes.io/hostname
      containers:
        - name: web
          image: nginxdemos/hello:plain-text
          ports:
            - containerPort: 80
---
apiVersion: v1
kind: Service
metadata:
  name: app-distribuida
spec:
  type: NodePort
  selector:
    app: distribuida
  ports:
    - port: 80
      targetPort: 80
      nodePort: 30080
\end{lstlisting}

Aplique o arquivo e veja onde os pods caíram:

\begin{lstlisting}[language=bash]
kubectl apply -f app-nuvem.yaml
kubectl get pods -o wide
\end{lstlisting}

\begin{tcolorbox}[colback=infobg, colframe=infoframe, arc=3pt, boxrule=0.8pt]
\textbf{OBSERVE:} Na coluna \texttt{NODE}, um pod deve estar no \texttt{lab-k8s-worker} e o outro no \texttt{lab-k8s-worker2}. O Kubernetes respeitou a regra de \textit{Anti-affinity} e não os colocou na mesma máquina. \textbf{Anote o nome exato do pod e do nó onde ele está rodando!}
\end{tcolorbox}

\subsection*{2. O Teste de Fogo (Derrubando um Servidor)}
Vamos provar que se um dos servidores cair, a aplicação sobrevive.

\textbf{1. Deixe o terminal observando os pods:}
\begin{lstlisting}[language=bash]
kubectl get pods -o wide -w
\end{lstlisting}

\textbf{2. Em outro terminal, isole um dos nós:} \textit{(Use o nome do worker onde um dos pods está)}
\begin{lstlisting}[language=bash]
kubectl cordon lab-k8s-worker
\end{lstlisting}

\textbf{3. Simule a queda do pod daquele nó:}
\begin{lstlisting}[language=bash]
kubectl delete pod <NOME_DO_POD_LOCAL>
\end{lstlisting}

\begin{tcolorbox}[colback=infobg, colframe=infoframe, arc=3pt, boxrule=0.8pt]
\textbf{O QUE ACONTECEU?} Olhe para o terminal do \textit{watch}. O pod que você deletou vai tentar renascer, mas vai ficar \textbf{\texttt{Pending}} (Pendente). O Kubernetes está protegendo a sua regra de \textit{Anti-affinity}: ele prefere deixar o pod pendente a colocar dois pods na mesma máquina sobrevivente.
\end{tcolorbox}

\subsection*{3. A Prova Final (O Terminal não mente)}
Deixe o \texttt{curl} rodando em loop para provar que o site não caiu:

\begin{lstlisting}[language=bash]
while ($true) { curl.exe -s http://localhost:30080 | Select-String "Server name"; Start-Sleep -Seconds 1 }
\end{lstlisting}

\begin{tcolorbox}[colback=infobg, colframe=infoframe, arc=3pt, boxrule=0.8pt]
\textbf{OBSERVE:} O terminal vai mostrar o nome do pod sobrevivente respondendo. O site \textbf{continua no ar}, servido 100\% pelo outro servidor! \textit{(Aperte \texttt{Ctrl+C} para parar).}
\end{tcolorbox}

\subsection*{4. Restaurando o Ambiente}
Para voltar tudo ao normal e o pod \texttt{Pending} voltar para o nó recuperado:

\begin{lstlisting}[language=bash]
kubectl uncordon lab-k8s-worker
\end{lstlisting}

\section*{Limpeza Final}

Para não sobrecarregar o seu computador, apague o cluster local:

\begin{lstlisting}[language=bash]
kubectl delete deployment app-distribuida
kubectl delete service app-distribuida
kind delete cluster --name lab-k8s
\end{lstlisting}

\end{document}